\chapter*{Abstract}
\addcontentsline{toc}{chapter}{Abstract}
\markboth{Abstract}{Abstract}

This project studies intrusion detection from IoT network flows and the review
of alerts in a monitoring application. It uses RT-IoT2022, checking its
123,117 rows and removing conflicting labels and duplicate observations.
Four model families are compared on training, validation, and test partitions
using three random seeds. The selection procedure chooses histogram-based
gradient boosting to distinguish attacks from normal traffic, followed by a
random forest to identify attack families.

A FastAPI service runs the models and stores observations, predictions, and
alerts. A persistent ingestion queue supports retries without creating duplicate
records. The React interface provides access to results, model versions, and
the history of analyst dispositions. Signature alerts from the Suricata
laboratory are identified by their source.

On the internal test split for seed~42, the cascade achieves a macro-F1 of
0.9520; the mean across the three seeds is 0.9686. Results for rare classes
remain fragile: the test contains only six \texttt{NMAP\_FIN\_SCAN} flows and
seven \texttt{Metasploit\_Brute\_Force\_SSH} flows. The study provides an internal
evaluation and an investigation prototype. A labelled NFStream collection and
evaluation on an external network are still needed to assess transfer.

\noindent\textbf{Keywords:} Internet of Things, intrusion detection, network
flows, machine learning, RT-IoT2022.
